% [IDEAL] The maximum principle, by counterexample. LEFT: the field the operator actually produces on
% a 5×5 grid toward a corner source — every interior node is strictly below its largest neighbour, so
% the via-arrows (steepest ascent) all escape to the source: no trap. RIGHT: a hand-authored FORBIDDEN
% configuration with a spurious interior maximum at the red node — every nearby arrow points INTO it
% and none escapes, a black hole. The screened operator can never produce this (Thm 5.4/5.6); the
% right panel is an impossibility, not a measurement. Data: maxp_legal.dat, maxp_forbidden.dat (figgen).
\providecommand{\figpath}{figures/}
\begin{tikzpicture}
\begin{axis}[
  name=L, width=6.4cm, height=6.4cm, axis equal image, enlargelimits=0.18, hide axis,
  colormap/viridis, point meta min=0, point meta max=1,
  title={\small legal: $f_i < \max_j f_j$ everywhere},
]
\addplot[point meta=explicit, quiver={u=\thisrow{dx}, v=\thisrow{dy}, scale arrows=0.34,
  every arrow/.append style={-{Stealth[length=4pt]}, black!65}}] table[meta=f] {\figpath maxp_legal.dat};
\addplot[scatter, only marks, scatter src=explicit, mark=*, mark size=3.4pt] table[meta=f] {\figpath maxp_legal.dat};
\node[circle, draw=black, fill=white, inner sep=1.3pt, label={[font=\tiny]below:src}] at (axis cs:0,0) {};
\end{axis}
\begin{axis}[
  name=R, at={(L.east)}, anchor=west, xshift=1.1cm,
  width=6.4cm, height=6.4cm, axis equal image, enlargelimits=0.18, hide axis,
  colormap/viridis, point meta min=0, point meta max=1,
  title={\small forbidden: an interior maximum traps flow},
]
\addplot[point meta=explicit, quiver={u=\thisrow{dx}, v=\thisrow{dy}, scale arrows=0.34,
  every arrow/.append style={-{Stealth[length=4pt]}, black!65}}] table[meta=f] {\figpath maxp_forbidden.dat};
\addplot[scatter, only marks, scatter src=explicit, mark=*, mark size=3.4pt] table[meta=f] {\figpath maxp_forbidden.dat};
\node[circle, draw=black, fill=white, inner sep=1.3pt, label={[font=\tiny]below:src}] at (axis cs:0,0) {};
\node[circle, draw=red, very thick, minimum size=13pt] at (axis cs:2,2) {};
\node[red, font=\tiny\bfseries, anchor=south] at (axis cs:2,2.55) {black hole};
\end{axis}
\end{tikzpicture}
